%%%%%%%%%%%%%%%%%%%%%%% CHAPTER - 6 %%%%%%%%%%%%%%%%%%%%\\
\chapter{System Implementation}
\label{C6} %%%%%%%%%%%%%%%%%%%%%%%%%%%%
\graphicspath{{Figures/Chapter-6figs/PDF/}{Figures/Chapter-6figs/}}
% \noindent\rule{\linewidth}{2pt}
%%%%%%%%%%%%%%%%%%%%%%%%%%%%%%%%%%%%%%%%%%%%%%%%%%%%%%%%%%%%%%%%%%%%%%%%%%%%%%%%%%
%%%%%%%%%%%%%%%%%%%%%%%%%%%%%%%%%%%%%%%%%%%%%%%% Section-1s........... %%%%%%%%%%%%%%%%%%%%%%%%%%%%%%%%%%%%%%%%%%%%%%%%%%%%%%%%%%%%%%%%%%%%%%%%%%%%
The System Implementation phase describes the practical realization of the proposed AI-Powered Cricket Batting Analysis System. This phase involves integrating deep learning, backend development, and mobile application development into a fully functional system. The implementation was carried out in a modular manner to ensure scalability, maintainability, and efficiency.
\section{Development Environment and Technology Stack} \label{S6.1}

The AI Cricket Coach system was developed using a combination of modern software tools and frameworks to ensure high performance, flexibility, and scalability. The backend was implemented using the Python programming language along with the Flask framework to build RESTful APIs for handling communication between the mobile application and the AI models. The core detection functionality was implemented using the YOLOv8 object detection model, while XGBoost regression was integrated to generate performance scores for each cricket shot. OpenCV was used for image preprocessing and manipulation tasks before model inference. The backend development and testing were carried out in Visual Studio Code. For model training, Google Colab with GPU acceleration was utilized to speed up training and optimize computational efficiency. The Ultralytics implementation of YOLOv8 and the XGBoost library were used for training and evaluation, while NumPy and Pandas were employed for feature processing and data handling. The mobile application was developed in Android Studio using Java, integrating CameraX for image capture, Retrofit for HTTP communication, Gson for JSON parsing, and the Text-to-Speech API for voice feedback. Proper dependency management and version compatibility were maintained to avoid runtime conflicts between modules.

\section{Dataset Implementation and Preparation} \label{S6.2}

The dataset forms the foundation of the AI Cricket Coach system, as model accuracy heavily depends on data quality. The dataset consists of over 8000 labeled cricket batting images categorized into Straight Drive, Pull Shot, Cut Shot, Sweep Shot, and an additional Unknown class to handle uncertain predictions. Each image was manually verified to ensure labeling correctness. For YOLOv8 training, bounding box annotation was performed to identify key regions such as the bat position, player body alignment, and shot execution frame. These annotations were stored in YOLO format using text files containing normalized coordinates. During preprocessing, all images were resized to a fixed resolution, typically 640×640 pixels, to match YOLOv8 input requirements. Pixel normalization was applied to improve model convergence during training. Data augmentation techniques such as horizontal flipping, rotation, brightness adjustment, and zoom scaling were used to enhance dataset diversity and reduce overfitting. The dataset was split into 70\% training data, 20\% testing data, and 10\% validation data to ensure reliable evaluation and generalization capability.

\section{YOLOv8 Model Implementation} \label{S6.3}

The system uses YOLOv8 for real-time cricket shot detection and classification instead of a traditional CNN classifier. YOLOv8 was selected due to its fast inference speed, real-time detection capability, high mean Average Precision (mAP), and lightweight architecture suitable for deployment. The architecture consists of a backbone network for feature extraction, a neck layer for feature aggregation, and a detection head responsible for bounding box prediction and classification. During training, pretrained weights were used as a base model and fine-tuned using the custom cricket dataset. The optimizer used was either SGD or Adam, depending on performance evaluation. The training process was conducted over multiple epochs to allow the model to learn important visual features such as bat angle, arm extension, body alignment, and follow-through direction. The loss function included bounding box loss, classification loss, and objectness loss to ensure accurate localization and classification. After successful training and validation, the model was exported as a .pt file for integration into the backend server.

\section{Feature Extraction and XGBoost} \label{S6.4}

To enhance the system beyond simple classification, a regression-based performance scoring mechanism was implemented using XGBoost. After YOLOv8 detects and classifies a cricket shot, numerical features are extracted from its output, including detection confidence score, bounding box dimensions, bat angle approximation, posture alignment, and shot stability indicators. These features are structured into numerical vectors and used to train the XGBoost regression model. XGBoost was chosen due to its high prediction accuracy, ability to handle non-linear relationships, and robustness against overfitting. The regression model predicts a performance score that reflects the quality and technical correctness of the shot. Evaluation metrics such as Root Mean Square Error (RMSE) and R-squared (R²) score were used to assess the model’s performance. The trained regression model was saved in .json or .pkl format and deployed within the Flask backend for real-time scoring.

\section{Backend Server Implementation} \label{S6.5}

The backend server functions as the core processing unit of the AI Cricket Coach system. It was developed using the Flask framework and provides a RESTful API endpoint named “/predict” to handle image prediction requests. When the Android application sends an image via HTTP POST request, the server receives the image, converts it into a NumPy array, and resizes it to match YOLOv8 input dimensions. The YOLOv8 model performs detection and classification, and relevant features are extracted from the detection results. These features are passed to the XGBoost regression model to compute a performance score. A confidence threshold mechanism is implemented to reduce misclassification; if the confidence score is below the predefined threshold, the shot is labeled as “Unknown.” The server then generates a structured JSON response containing the predicted shot type, confidence value, performance score, and feedback message, which is sent back to the Android application.

\section{Android Application Implementation} \label{S6.6}

The Android application serves as the user interface of the system and allows players to interact with the AI Cricket Coach. It was developed in Android Studio using Java. The CameraX API was integrated to enable real-time image capture of batting actions. After capturing an image, the application processes and prepares it for transmission to the backend server. Retrofit was used to send the image as a multipart HTTP request and handle asynchronous network communication. Upon receiving the JSON response, the application parses the data using Gson and dynamically updates the user interface to display the predicted shot type and performance score. The Text-to-Speech API was integrated to provide instant voice feedback, enhancing user experience. Logout and session management features were implemented to maintain secure user access. Error handling mechanisms were added to manage network timeouts, server unavailability, and invalid responses to ensure reliability.

\section{Client–Server Integration} \label{S6.7}

The AI Cricket Coach system follows a RESTful client–server architecture to ensure efficient communication between the Android application and the backend server. When a user captures an image, the Android application converts it into a multipart HTTP request and sends it to the Flask server. The server processes the image using YOLOv8 for detection and XGBoost for regression scoring, and then returns a structured JSON response. The Android application receives this response, displays the results on the screen, and announces them through voice feedback. This separation of responsibilities ensures that computationally intensive tasks are handled by the server, reducing processing load on the mobile device and improving overall system performance.

\section{Testing and Validation} \label{S6.8}

Comprehensive testing and validation were conducted to ensure system accuracy and reliability. The YOLOv8 model was evaluated using the validation dataset, and detection performance was measured using metrics such as mean Average Precision (mAP). The XGBoost regression model was assessed using statistical evaluation measures including RMSE and R-squared score. API testing was performed using Postman to validate request handling, response structure, and latency. Integration testing ensured smooth communication between the Android application and the backend server. Real-world testing was also conducted using live cricket practice images to verify shot classification accuracy and consistency of performance scoring under practical conditions.

\section{Deployment and Execution} \label{S6.9}

The Flask backend server was deployed locally and executed using Python commands in Visual Studio Code. The Android application was configured to connect to the backend using the server’s IP address within the same network environment. This setup enabled real-time communication between the client and server. The modular system architecture allows for future deployment on cloud platforms such as AWS, Render, or Heroku, enabling public API access and multi-user scalability.
